\documentclass[12pt]{article}
\usepackage[utf8]{inputenc}
\usepackage[T1]{fontenc}
\usepackage[margin=1in]{geometry}
\usepackage{setspace}
\usepackage{xcolor}
\usepackage[colorlinks=true, urlcolor=blue!60!black, linkcolor=blue!60!black]{hyperref}
\setlength{\parindent}{0em}
\setlength{\parskip}{0.6em}
\renewcommand{\baselinestretch}{1.0}
\begin{document}
\begin{center}
{\LARGE \textit{Citizens United v. Federal Election Commission}, 558 U.S. 310 (2010)\par}
\vspace{0.5em}
{\large Case Brief --- Two-Page Edition\par}
\end{center}

\section{Memory Jogger}
The government may not ban corporate or union independent political expenditures (5--4), but it may require disclaimers and disclosure (8--1) (\href{https://www.courtlistener.com/opinion/1741/citizens-united-v-federal-election-commission/}{CourtListener}).

\section{Detailed Case Facts}
Citizens United, a nonprofit corporation, released \textit{Hillary: The Movie}, a documentary attacking Senator Hillary Clinton during the 2008 primaries, and planned to offer it on video-on-demand and promote it with television ads. Section 441b, as amended by McCain--Feingold, barred corporations from using treasury funds for ``electioneering communications'' naming a candidate within 30 days of a primary or 60 days of a general election; corporations could speak only through PACs.

\section{Procedural History}
A three-judge district court granted the FEC summary judgment. On direct appeal, the Supreme Court ordered reargument on whether to overrule \textit{Austin} and \textit{McConnell}.

\section{Judicial Votes}
Kennedy wrote for Roberts, Scalia, Alito, and Thomas (5--4); Stevens, Ginsburg, Breyer, and Sotomayor joined Part IV on disclosure (8--1). Roberts and Scalia concurred; Stevens and Thomas each concurred in part and dissented in part.

\section{Holding}
The government may not suppress political speech on the basis of the speaker's corporate identity. Independent expenditures do not give rise to \textit{quid pro quo} corruption, so the ban fails strict scrutiny; \textit{Austin} is overruled. Disclosure survives exacting scrutiny.

\section{Analysis of Opinions}
\textbf{Majority:} Section 441b was an outright ban, the PAC option no cure; government has no interest in equalizing voices, and ingratiation and access are not corruption. \textbf{Roberts:} ``There is a difference between judicial restraint and judicial abdication.'' \textbf{Scalia:} the First Amendment protects speech, not categories of speakers. \textbf{Stevens:} the Court reached an abandoned facial challenge; corporations cannot vote or run for office, and the Framers saw corruption more broadly. \textbf{Thomas:} disclosure exposed Proposition 8 donors to reprisals.

\section{Examples: Future Applications}
\textbf{Same-Side:} a state ban on corporate independent expenditures falls (\textit{American Tradition Partnership v. Bullock}). \textbf{Opposite-Side:} the ban on direct corporate contributions to candidates survives. \textbf{Fence-Sitter:} a ban on spending by foreign-owned U.S. subsidiaries, a question the Court left open.

\section{Critique}
Critics such as Richard Hasen say Super PACs and ``dark money'' refute the claim that independent spending does not corrupt, and Zephyr Teachout argues the Framers' idea of corruption was broader. Defenders such as Joel Gora say the law silenced small businesses and nonprofits and could have banned books.

\section{Key Quotations}
\begin{quote}
``For the reasons explained above, we now conclude that independent expenditures, including those made by corporations, do not give rise to corruption or the appearance of corruption.''
\end{quote}
\noindent --- Justice Kennedy, Opinion of the Court, 558 U.S. at 357.

\end{document}
